% GENERATED FILE. Edit canonical lessons, then run npm run generate:books.
% canonical-source-hash: fnv1a64:80b97b9c116a136d
% canonical-lessons: GU-C231-goli, GU-C231-malam, GU-C231-nars, GU-C231-sarvar, GU-C231-matvu

\chapter{Pills and ointment}
\label{ch:gu-pills-and-ointment}
\hlchaptermodality{231}

\begin{chapteropening}
\textbf{By the end of this chapter:} \emph{I can ask about pills, ointment and treatment.}

Complete the last lesson of chapter 231: I can ask about pills, ointment and treatment.
\end{chapteropening}

\section[\texorpdfstring{\gu{ગોળી} (goḷī)}{ગોળી (goḷī)}]{\gu{ગોળી} (goḷī) — pill, tablet}

\label{lesson:GU-C231-goli}



\begin{quote}
Before the new one: say the Gujarati for to suffer, then the Gujarati for to faint.
\end{quote}

\subsection*{You'll want to know: \gu{ગોળી}}
\textbf{\gu{ગોળી}} — \emph{goḷī} — \textquotedblleft{}pill, tablet\textquotedblright{}.

\begin{quote}
\textbf{\gu{હું} \gu{રાત્રે} \gu{ગોળી} \gu{નથી} \gu{લેતો}.} — \emph{hũ rātre goḷī nathī leto} — \textquotedblleft{}I don't take the pill at night.\textquotedblright{}
\end{quote}

A feminine noun.

\begin{cousinweb}
\textbf{\gu{ગોળી}} (\emph{goḷī}) is a tadbhava. It comes from Sanskrit \emph{gola}, \textquotedblleft{}ball\textquotedblright{}, through \emph{golikā}, \textquotedblleft{}little ball\textquotedblright{}, and CDIAL glosses the Gujarati word \textquotedblleft{}pill\textquotedblright{} (CDIAL 4321).
\end{cousinweb}

\begin{grammarlens}[title={What you've built}]
One more word for everyday life at A2.
\end{grammarlens}

\subsection*{Your turn}
\begin{itemize}
\raggedright
  \item \emph{Say it:} \emph{goḷī}
  \item \emph{Say it:} \emph{goḷī}, once more
  \item \emph{Say it:} \emph{hũ rātre goḷī nathī leto} — I don't take the pill at night.
  \item \emph{Recall:} say the Gujarati for to suffer, then the Gujarati for to faint, then say \emph{goḷī} again
\end{itemize}

\begin{tcolorbox}[breakable,colback=teal!4,colframe=teal!35!black,title={Before you move on}]
What does \gu{ગોળી} mean? (\textquotedblleft{}Pill, tablet\textquotedblright{}.) Say it once more.
\end{tcolorbox}
\section[\texorpdfstring{\gu{મલમ} (malam)}{મલમ (malam)}]{\gu{મલમ} (malam) — ointment, balm}

\label{lesson:GU-C231-malam}



\begin{quote}
Before the new one: say the Gujarati for to faint, then the Gujarati for pill.
\end{quote}

\subsection*{You'll want to know: \gu{મલમ}}
\textbf{\gu{મલમ}} — \emph{malam} — \textquotedblleft{}ointment, balm\textquotedblright{}.

\begin{quote}
\textbf{\gu{શું} \gu{આ} \gu{મલમ} \gu{સારો} \gu{છે}?} — \emph{śũ ā malam sāro che?} — \textquotedblleft{}Is this ointment good?\textquotedblright{}
\end{quote}

\begin{cousinweb}
\textbf{\gu{મલમ}} (\emph{malam}) is a loan, usually traced to Arabic \emph{marham}, \textquotedblleft{}ointment\textquotedblright{}, which reached Gujarati through Persian. Its Arabic root is \emph{r-h-m} (Lane).
\end{cousinweb}

\begin{grammarlens}[title={What you've built}]
One more word for everyday life at A2.
\end{grammarlens}

\subsection*{Your turn}
\begin{itemize}
\raggedright
  \item \emph{Say it:} \emph{malam}
  \item \emph{Say it:} \emph{malam}, once more
  \item \emph{Say it:} \emph{śũ ā malam sāro che?} — Is this ointment good?
  \item \emph{Recall:} say the Gujarati for to faint, then the Gujarati for pill, then say \emph{malam} again
\end{itemize}

\begin{tcolorbox}[breakable,colback=teal!4,colframe=teal!35!black,title={Before you move on}]
What does \gu{મલમ} mean? (\textquotedblleft{}Ointment, balm\textquotedblright{}.) Say it once more.
\end{tcolorbox}
\section[\texorpdfstring{\gu{નર્સ} (nars)}{નર્સ (nars)}]{\gu{નર્સ} (nars) — nurse}

\label{lesson:GU-C231-nars}



\begin{quote}
Before the new one: say the Gujarati for pill, then the Gujarati for ointment.
\end{quote}

\subsection*{You'll want to know: \gu{નર્સ}}
\textbf{\gu{નર્સ}} — \emph{nars} — \textquotedblleft{}nurse\textquotedblright{}.

\begin{quote}
\textbf{\gu{શું} \gu{નર્સ} \gu{દવાખાનામાં} \gu{છે}?} — \emph{śũ nars davākhānāmā̃ che?} — \textquotedblleft{}Is the nurse at the clinic?\textquotedblright{}
\end{quote}

Mostly feminine.

\begin{cousinweb}
\textbf{\gu{નર્સ}} (\emph{nars}) is a loan from English \emph{nurse}.
\end{cousinweb}

\begin{grammarlens}[title={What you've built}]
One more word for everyday life at A2.
\end{grammarlens}

\subsection*{Your turn}
\begin{itemize}
\raggedright
  \item \emph{Say it:} \emph{nars}
  \item \emph{Say it:} \emph{nars}, once more
  \item \emph{Say it:} \emph{śũ nars davākhānāmā̃ che?} — Is the nurse at the clinic?
  \item \emph{Recall:} say the Gujarati for pill, then the Gujarati for ointment, then say \emph{nars} again
\end{itemize}

\begin{tcolorbox}[breakable,colback=teal!4,colframe=teal!35!black,title={Before you move on}]
What does \gu{નર્સ} mean? (\textquotedblleft{}Nurse\textquotedblright{}.) Say it once more.
\end{tcolorbox}
\section[\texorpdfstring{\gu{સારવાર} (sārvār)}{સારવાર (sārvār)}]{\gu{સારવાર} (sārvār) — treatment, medical care}

\label{lesson:GU-C231-sarvar}



\begin{quote}
Before the new one: say the Gujarati for ointment, then the Gujarati for nurse.
\end{quote}

\subsection*{You'll want to know: \gu{સારવાર}}
\textbf{\gu{સારવાર}} — \emph{sārvār} — \textquotedblleft{}treatment, medical care\textquotedblright{}.

\begin{quote}
\textbf{\gu{શું} \gu{અહીં} \gu{સારવાર} \gu{મોંઘી} \gu{છે}?} — \emph{śũ ahī̃ sārvār mõghī che?} — \textquotedblleft{}Is treatment expensive here?\textquotedblright{}
\end{quote}

\begin{cousinweb}
The origin of \textbf{\gu{સારવાર}} (\emph{sārvār}) is not settled: CDIAL lists no Gujarati form for it, and its make-up is not clear.
\end{cousinweb}

\begin{grammarlens}[title={What you've built}]
One more word for everyday life at A2.
\end{grammarlens}

\subsection*{Your turn}
\begin{itemize}
\raggedright
  \item \emph{Say it:} \emph{sārvār}
  \item \emph{Say it:} \emph{sārvār}, once more
  \item \emph{Say it:} \emph{śũ ahī̃ sārvār mõghī che?} — Is treatment expensive here?
  \item \emph{Recall:} say the Gujarati for ointment, then the Gujarati for nurse, then say \emph{sārvār} again
\end{itemize}

\begin{tcolorbox}[breakable,colback=teal!4,colframe=teal!35!black,title={Before you move on}]
What does \gu{સારવાર} mean? (\textquotedblleft{}Treatment, medical care\textquotedblright{}.) Say it once more.
\end{tcolorbox}
\section[\texorpdfstring{\gu{મટવું} (maṭvũ)}{મટવું (maṭvũ)}]{\gu{મટવું} (maṭvũ) — to be cured, to clear up}

\label{lesson:GU-C231-matvu}



\begin{quote}
Before the new one: say the Gujarati for nurse, then the Gujarati for treatment.
\end{quote}

\subsection*{You'll want to know: \gu{મટવું}}
\textbf{\gu{મટવું}} — \emph{maṭvũ} — \textquotedblleft{}to be cured, to clear up\textquotedblright{}.

\begin{quote}
\textbf{\gu{શું} \gu{તમારો} \gu{તાવ} \gu{મટ્યો}?} — \emph{śũ tamāro tāv maṭyo?} — \textquotedblleft{}Has your fever gone?\textquotedblright{}
\end{quote}

\begin{cousinweb}
\textbf{\gu{મટવું}} (\emph{maṭvũ}) is probably a tadbhava. CDIAL lists it, glossed \textquotedblleft{}to be rubbed out\textquotedblright{}, under Sanskrit \emph{mṛṣṭa}, \textquotedblleft{}rubbed, washed\textquotedblright{}, though it calls one change of sound unexplained (CDIAL 10299); an illness that is cured is rubbed out.
\end{cousinweb}

\begin{grammarlens}[title={What you've built}]
One more word for everyday life at A2.
\end{grammarlens}

\subsection*{Your turn}
\begin{itemize}
\raggedright
  \item \emph{Say it:} \emph{maṭvũ}
  \item \emph{Say it:} \emph{maṭvũ}, once more
  \item \emph{Say it:} \emph{śũ tamāro tāv maṭyo?} — Has your fever gone?
  \item \emph{Recall:} say the Gujarati for nurse, then the Gujarati for treatment, then say \emph{maṭvũ} again
\end{itemize}

\begin{tcolorbox}[breakable,colback=teal!4,colframe=teal!35!black,title={Before you move on}]
What does \gu{મટવું} mean? (\textquotedblleft{}To be cured, to clear up\textquotedblright{}.) Say it once more.
\end{tcolorbox}
